\documentclass[11pt]{article}

\usepackage{amsmath,amssymb}
\usepackage{graphicx}
\usepackage{float}
\usepackage{pgf,tikz,pgfplots}
\usepackage{mathrsfs}
\usepackage{color}
\usepackage{soulutf8}
\usepackage{siunitx}
\usepackage{gensymb}
\usepackage{textcomp}
\usepackage[font=footnotesize, justification=justified]{caption}
\usepackage[font=scriptsize]{subcaption}
\usepackage{listingsutf8}
\usepackage{multirow}
\usepackage{makecell}
\usepackage{ragged2e}
\usepackage{tabularx}
\usepackage{url}
\usepackage{adjustbox}
\usepackage{animate}
\usepackage{etoolbox}
\usepackage{lmodern}
\usepackage{booktabs}
\usepackage{bm}
\usepackage{algorithm}
\usepackage{algpseudocode}
\usepackage{amsthm,thmtools}
\usepackage{scalerel}
\usepackage{lmodern}
\usepackage{mdframed}
\usepackage{multicol}
\usepackage{graphicx} % Required for inserting 
\usepackage{listings}
\usepackage{enumitem}
\usepackage[T1]{fontenc}
\usepackage[brazil]{babel}
\usepackage{setspace}
\usepackage{titlesec}
\usepackage{hyperref}

\usepackage{hyperref}
\hypersetup{
    colorlinks=true,
    linkcolor=blue,
    urlcolor=blue
}



\usepackage{geometry}
\geometry{a4paper,
        width=150mm, %ok
            top=20mm, %ok
            bottom=40mm,
            bindingoffset=3mm
        }
\usepackage{fancyhdr}
\pagestyle{fancy}
\fancyhead{}
\fancyhead[L]{\noindent{\Large \textbf{Ciência de Dados e Inteligência Artificial}}\\
                {\normalsize \textbf{Disciplina:} Computação Escalável} \\
                {\normalsize \textbf{Professor:} Larissa C. Shimomura}}
\fancyhead[R]{\includegraphics[width = 0.4\textwidth]{logo_uel.png}
}
\renewcommand{\headrulewidth}{0pt}
\setlength{\headsep}{2.5cm}
\fancyheadoffset{1cm}


%\definecolor{numbers}{HTML}{9E9E9E}
\definecolor{keywords}{HTML}{E91E63}
\definecolor{comments}{HTML}{9E9E9E}
\definecolor{strings}{HTML}{AFB42B}
\definecolor{environments}{HTML}{03A9F4}
\definecolor{options}{HTML}{FF5722}

\definecolor{link_color}{RGB}{51,51,178}

\lstset {
	basicstyle=\fontfamily{cmtt}\selectfont,
	%numbers=left,
	%numberstyle=\tiny\color{numbers},
	%stepnumber=1,
	%numbersep=10pt,
	backgroundcolor=\color{gray!20},
	showspaces=false,
	showstringspaces=false,
	showtabs=false,
	frame=single,
	rulecolor=\color{black},
	tabsize=3,
	captionpos=b,
	breaklines=true,
	breakatwhitespace=false,
	% keywordstyle=\color{keywords},
	% keywordstyle=[2]\color{environments},
	% keywordstyle=[3]\color{options},
	% commentstyle=\color{comments},
	% stringstyle=\color{strings},
	upquote=true,
	inputencoding=utf8/latin1,
	extendedchars=true,
    % keepspaces=true,
	% identifierstyle=\color{black},
	columns=fullflexible,
        literate=
        {á}{{\'a}}1
        {à}{{\`a}}1
        {ã}{{\~a}}1
        {é}{{\'e}}1
        {ê}{{\^e}}1
        {í}{{\'i}}1
        {ó}{{\'o}}1
        {õ}{{\~o}}1
        {ú}{{\'u}}1
        {ü}{{\"u}}1
        {ç}{{\c{c}}}1
}


\setstretch{1.2}
\titleformat{\section}{\normalfont\bfseries\uppercase}{\thesection.}{1em}{}
\titleformat{\subsection}{\normalfont\bfseries}{\thesubsection}{1em}{}

\begin{document}

\begin{center}
    \Large\textbf{Relatório Final: Projeto Prático de Processamento Paralelo e Distribuído}\\
    \textbf{Integrantes: Aluno A, Aluno B, Aluno C, Aluno D}
\end{center}


\section{Introdução}
% Escreva aqui uma breve introdução sobre a tarefa de Ciência de Dados escolhida e uma visão geral dos dois algoritmos selecionados para o estudo comparativo.

\section{Algoritmos Implementados}

\subsection{Algoritmo 1: [Nome do Algoritmo] (Responsáveis: Nome do Aluno A e Aluno B)}
% Detalhe o algoritmo, seus passos e as escolhas de projeto do primeiro algoritmo. Explique explicitamente as decisões tomadas para implementar o paralelismo ou a distribuição (ex: uso de Threads, ExecutorService, ForkJoin ou Spark). 

\subsection{Algoritmo 2: [Nome do Algoritmo] (Responsáveis: Nome do Aluno C e Aluno D)}
% Detalhe o algoritmo, seus passos e as escolhas de projeto do primeiro algoritmo. Explique explicitamente as decisões tomadas para implementar o paralelismo ou a distribuição (ex: uso de Threads, ExecutorService, ForkJoin ou Spark). 

\section{Setup Experimental}

\subsection{Datasets Escolhidos}
% Descreva os dois datasets utilizados nos testes (origem dos dados, se são reais ou sintéticos, características e distribuições). Lembre-se que o tamanho deve variar no intervalo recomendado (de 10.000 a 150.000 pontos, com aumentos graduais, ex: passo de 20k em 20k).

\subsection{Linguagem e Configuração do Ambiente}
% Informe a linguagem utilizada (Java ou Scala/Spark) e descreva o ambiente de hardware e software utilizado para os testes (número de núcleos/cores da CPU, limite de memória alocado para a JVM). 

\section{Análise de Escalabilidade}
% Apresente os gráficos comparativos (Tamanho do Dataset vs. Tempo de Execução) para cada dataset. Analise o comportamento das duas linhas/séries correspondentes aos algoritmos. Discuta os gargalos observados e relate se os resultados condizem com a teoria ou com o esperado para a arquitetura paralela escolhida.

\section{Conclusões}
% Resuma as principais conclusões obtidas a partir dos experimentos práticos. Comente qual dos algoritmos apresentou melhor desempenho e escalabilidade sob diferentes volumes de dados. Aponte limitações da implementação atual e sugira possíveis melhorias ou trabalhos futuros.

\section*{Declaração de Uso de IA}
% Declare explicitamente se a Inteligência Artificial foi utilizada no projeto (conforme as regras da disciplina: permitida apenas para leitura inicial de arquivos, scripts de automação de testes e plots de resultados). 

\end{document}